\begin{helper}[Shifted closing-cycle coloring]
Let $N$ be balanced normal globally maximum data for
$(G,k,e_0,\varphi,T)$, and fix
$j\in\operatorname{Fin}(N.height)$ with $1\leq j$.  There are a partial
coloring $\psi$ and distinct colors $A,B$ such that:
\begin{enumerate}
\item if $j$ is odd, $(A,B)=(N.beta,N.alpha)$, while if $j$ is even,
      $(A,B)=(N.alpha,N.beta)$;
\item $\psi$ is proper and leaves exactly
      $f_j=\noderef{NormalCycleEdge}(N,j)$ uncolored;
\item $A$ is missing at $y_j$, and $B$ is missing at
      $y_{(j+1)\bmod N.height}$;
\item for every $0\leq i<j$, $\psi(f_i)=\varphi(f_{i+1})$;
\item every exact edge identity outside $\{f_0,\ldots,f_j\}$ retains its
      old value; and
\item for every old tree-edge position $r\geq N.height-1$,
      $\psi(T.edge(r))=\varphi(T.edge(r))$.
\end{enumerate}
All bounds are the corresponding finite-index bounds in the Lean
statement.
\end{helper}
\begin{proof}
Write $p=N.height$ and use the exact identities and alternating colors
from \noderef{ScheideTrunkRotationCycle}.  Let $B$ be the old color of
$f_j$ and let $A$ be the other one: explicitly, if $j$ is odd take
$(A,B)=(N.beta,N.alpha)$, and if $j$ is even take
$(A,B)=(N.alpha,N.beta)$.  The cited alternation proves that this is the
old color of $f_j$, and its proof also gives $N.alpha\ne N.beta$; hence
$A\ne B$.

Define $\psi$ on exact edge identities by
\[
 \psi(f_j)=\operatorname{none},\qquad
 \psi(f_i)=\varphi(f_{i+1})\ (0\leq i<j),\qquad
 \psi(g)=\varphi(g)\quad
   (g\notin\{f_0,\ldots,f_j\}).                         \tag{1}
\]
Injectivity of $i\mapsto f_i$ makes the cases disjoint and hence defines a
function.  Each $f_i$ with $i<j$ receives a colored value by the
alternation theorem.  If $g\ne f_j$ is not among the earlier shifted
identities, then either it is one of the colored cycle identities or it is
unchanged and differs from $f_0=e_0$.  The old unique-uncolored clause
\noderef{IsOneEdgeUncolored} then supplies a color.  Thus exactly $f_j$
is uncolored.

We verify properness locally.  For $0\leq i<j$, let $c_i$ be the old
color of $f_{i+1}$, hence the new color of $f_i$.  At $y_{i+1}$, old
properness makes $f_{i+1}$ the unique incident carrier of $c_i$: any
distinct second carrier would be adjacent by \noderef{AdjacentEdges} and
contradict \noderef{IsProperEdgeColoring}.  Under (1), that unique old
carrier is recolored if $i+1<j$ and becomes uncolored if $i+1=j$.  Hence
no unchanged edge at this endpoint conflicts with the new color on
$f_i$.

At the other endpoint $y_i$, there are three cases.  For $i=0$,
$c_i=N.alpha$, which is missing at $y_0$.  For $i=1$,
$c_i=N.beta$, which is missing at $y_1$.  In either case
\noderef{MissingColors} and \noderef{PresentColors} exclude an unchanged
carrier.  For $2\leq i<j$, alternation gives
$\varphi(f_{i-1})=c_i$.  Old properness makes $f_{i-1}$ the unique old
carrier at $y_i$, and (1) recolors it simultaneously with the other trunk
color.  Again no unchanged identity conflicts.

Two cycle identities share an endpoint only when their indices are
consecutive modulo $p$, by the exact endpoint formula, vertex injectivity
in \noderef{TreeSequence}, and $p\geq3$.  Consecutive shifted identities
receive consecutive old colors and are different.  At $y_0$, the
additional comparison is between $f_0$, which receives $N.alpha$, and
$f_{p-1}$, which remains $N.beta$ unless it is the new uncolored identity.
At the other end $f_j$ is uncolored.  All pairs not involving a shifted
identity retain their proper old values, and the preceding endpoint
analysis covers every mixed pair.  This proves properness of $\psi$.

At $y_j$, if $j=1$, $A=N.beta$ was already missing and the only incoming
new color is $B=N.alpha$.  If $j>1$, alternation makes $f_{j-1}$ the
unique old $A$-carrier there; it is recolored to $B$, while $f_j$ becomes
uncolored.  No edge changes to $A$, so $A$ is missing at $y_j$.  At
$y_{(j+1)\bmod p}$, old properness makes $f_j$ the unique old $B$-carrier,
and it becomes uncolored.  The other incident cycle identity has new color
$A$: it is the unchanged $f_{j+1}$ if $j<p-1$, and is the newly colored
$f_0$ if $j=p-1$.  No noncycle identity changes.  Thus $B$ is missing at
that successor vertex.  This proves the two exact missing-color clauses.

The fourth and fifth conclusions are the corresponding clauses of (1).
For the sixth, final tail clause, let $r\geq p-1$.  Each $f_i$ with
$0\leq i\leq p-2$ is the old tree edge at position $i$, so old edge
injectivity separates it from $T.edge(r)$.  If the closing identity
$f_{p-1}$ equaled $T.edge(r)$, the attachment field of
\noderef{TreeSequence} would make it incident with the fresh vertex
$y_{r+1}$, while its endpoint formula makes both ends $y_{p-1},y_0$.
Vertex injectivity and $r+1\geq p$ contradict this.  Therefore the tail
identity is outside every shifted $f_i$, and the fixed clause in (1) gives
the asserted equality of its values.
\end{proof}
